\section{Related Work}
\label{sec:related}

\para{Exaggerated safety and over-refusal.}
The closest line of work to ours studies \emph{over-refusal}---models refusing safe prompts due to overgeneralized safety training. \citet{rottger2023xstest} introduced \xstest, a benchmark of 250~safe and 200~unsafe prompts designed to detect exaggerated safety behaviors. They found that models like Llama~2 refuse over 38\% of safe prompts due to lexical overfitting: the word ``killing'' triggers refusal regardless of whether the context is violent or benign (\eg ``killing time''). \citet{cui2024orbench} scaled this analysis with \orbench, an 80K-prompt benchmark evaluated across 32~LLMs, finding a Spearman correlation of 0.89 between safety and over-refusal---suggesting the two are nearly inseparable with current training methods. \citet{xie2024sorrybench} contributed a fine-grained 45-category taxonomy for evaluating refusal behaviors. Unlike all of these works, we do not measure binary refusal but instead measure the \emph{quality} of engagement when models do respond, revealing a richer and previously unmeasured space of avoidance behavior.

\para{Refusal mechanisms and intervention.}
Recent mechanistic work has investigated \emph{how} refusal is implemented in model internals. \citet{arditi2024surgical} showed that refusal can be removed via ablation of a single direction in activation space, suggesting refusal is encoded as a simple linear feature. Subsequent work has complicated this picture: refusal behavior has been shown to involve nonlinear components and to be consistent across languages \cite{arditi2024surgical}. These mechanistic insights focus on explicit refusal circuits; whether implicit avoidance involves similar or distinct mechanisms remains an open question that our behavioral measurements motivate.

\para{Sycophancy and RLHF artifacts.}
Sycophancy---excessive agreement with user views---represents another class of RLHF-induced behavioral artifacts. \citet{sharma2023sycophancy} demonstrated that sycophancy is a general behavior across tasks and that it \emph{increases} with more RLHF training, implicating the training process itself. \citet{malmqvist2024sycophancy} surveyed causes and mitigations, framing sycophancy and over-refusal as ``two sides of the same coin''---both emerge from human preference optimization that rewards pleasing responses. Our work extends this perspective: we argue that implicit avoidance is a third artifact of this same process, where models learn to hedge and qualify on topics that training data associates with social sensitivity.

\para{Identity-related speech suppression.}
\citet{proebsting2025speech} audited five content moderation APIs across seven datasets and nine identity groups, finding that identity-related speech is disproportionately flagged as harmful. Critically, suppression patterns follow \emph{stereotypical associations} from training data: disability-related speech triggers self-harm flags, LGBT-related speech triggers sexual content flags. This finding directly supports our hypothesis that avoidance patterns generalize from training data rather than from explicit safety rules.

\para{Safety--helpfulness tradeoff.}
Broader work on the tension between safety and helpfulness provides context for our findings. \citet{bai2022constitutional} introduced constitutional AI, training models to be harmless using AI feedback. \citet{ouyang2022instructgpt} established RLHF methodology for instruction following. Both approaches create models that must balance helpfulness against caution, and our results suggest this balance is not uniform across topics---models are systematically more cautious on some topics than others, even when all topics are safe to discuss.

\para{Position of our work.}
While prior work has measured explicit refusal (binary: refused or not), we introduce the concept of \emph{implicit avoidance}---a continuous, multi-signal phenomenon where models engage with topics but with measurably degraded directness. Our cross-model convergence analysis further distinguishes our contribution: by showing that avoidance patterns correlate strongly across model sizes, we provide evidence about the \emph{origin} of these patterns, not just their existence.
